Practitioners often aim to infer an unobserved population trajectory using sample snapshots at multiple time points. E.g. given single-cell sequencing data, scientists would like to learn how gene expression changes over a cell’s life cycle. But sequencing any cell destroys that cell. 
So we can access data for any particular cell only at a single time point, but we have data across many cells. The deep learning community has recently explored using Schrödinger bridges (SBs) and their extensions in similar settings. However, existing methods either (1) interpolate between just two time points or (2) require a single fixed reference dynamic (often set to Brownian motion within SBs). But learning piecewise from adjacent time points can fail to capture long-term dependencies. And practitioners are typically able to specify a model family for the reference dynamic but not the exact values of the parameters within it. So we propose a new method that (1) learns the unobserved trajectories from sample snapshots across multiple time points and (2) requires specification only of a family of reference dynamics, not a single fixed one. We demonstrate the advantages of our method on simulated and real data. %\revised{Code to reproduce the experiments is available at \href{https://github.com/YunyiShen/SB-Iterative-Reference-Refinement}{https://github.com/YunyiShen/SB-Iterative-Reference-Refinement}.}